\section{AHSME 1995}
        \begin{enumerate}
        	\item (\textit{AHSME 1995}) Kim earned scores of 87, 83, and 88 on her first three mathematics examinations. If Kim receives a score of 90 on the fourth exam, then her average will 


 $\mathrm{(A) \ \text{remain the same} } \qquad \mathrm{(B) \ \text{increase by 1} } \qquad \mathrm{(C) \ \text{increase by 2} } \qquad \mathrm{(D) \ \text{increase by 3} } \qquad \mathrm{(E) \ \text{increase by 4} }$


 

	\item (\textit{AHSME 1995}) If $\sqrt {2 + \sqrt {x}} = 3$, then $x =$


 $\mathrm{(A) \ 1 } \qquad \mathrm{(B) \ \sqrt{7} } \qquad \mathrm{(C) \ 7 } \qquad \mathrm{(D) \ 49 } \qquad \mathrm{(E) \ 121 }$


 

	\item (\textit{AHSME 1995}) The total in-store price for an appliance is $\textdollar 99.99$. A television commercial advertises the same product for three easy payments of $\textdollar 29.98$ and a one-time shipping and handling charge of $\textdollar 9.98$. How many cents are saved by buying the appliance from the television advertiser?


 $\mathrm{(A) \ 6 } \qquad \mathrm{(B) \ 7 } \qquad \mathrm{(C) \ 8 } \qquad \mathrm{(D) \ 9 } \qquad \mathrm{(E) \ 10 }$


 

	\item (\textit{AHSME 1995}) If $M$ is $30 \%$ of $Q$, $Q$ is $20 \%$ of $P$, and $N$ is $50 \%$ of $P$, then $\frac {M}{N} =$


 $\mathrm{(A) \ \frac {3}{250} } \qquad \mathrm{(B) \ \frac {3}{25} } \qquad \mathrm{(C) \ 1 } \qquad \mathrm{(D) \ \frac {6}{5} } \qquad \mathrm{(E) \ \frac {4}{3} }$


 

	\item (\textit{AHSME 1995}) A rectangular field is 300 feet wide and 400 feet long. Random sampling indicates that there are, on the average, three ants per square inch through out the field. [12 inches = 1 foot.] Of the following, the number that most closely approximates the number of ants in the field is 


 $\mathrm{(A) \ \text{500 thousand} } \qquad \mathrm{(B) \ \text{5 million} } \qquad \mathrm{(C) \ \text{50 million} } \qquad \mathrm{(D) \ \text{500 million} } \qquad \mathrm{(E) \ \text{5 billion} }$


 

	\item (\textit{AHSME 1995}) The figure shown can be folded into the shape of a cube. In the resulting cube, which of the lettered faces is opposite the face marked $x$? 


 
                    \begin{figure}[H]
                        \centering
                        \includegraphics[width=0.5\textwidth]{Problem/AHSME_I_1995_6_Problem_Figure_1.png}
                    \end{figure}
                    


 $\mathrm{(A) \ A } \qquad \mathrm{(B) \ B } \qquad \mathrm{(C) \ C } \qquad \mathrm{(D) \ D } \qquad \mathrm{(E) \ E }$


 

	\item (\textit{AHSME 1995}) The radius of Earth at the equator is approximately 4000 miles. Suppose a jet flies once around Earth at a speed of 500 miles per hour relative to Earth. If the flight path is a negligible height above the equator, then, among the following choices, the best estimate of the number of hours of flight is: 


 $\mathrm{(A) \ 8 } \qquad \mathrm{(B) \ 25 } \qquad \mathrm{(C) \ 50 } \qquad \mathrm{(D) \ 75 } \qquad \mathrm{(E) \ 100 }$


 

	\item (\textit{AHSME 1995}) In $\triangle ABC$, $\angle C = 90^\circ, AC = 6$ and $BC = 8$. Points $D$ and $E$ are on $\overline{AB}$ and $\overline{BC}$, respectively, and $\angle BED = 90^\circ$. If $DE = 4$, then $BD =$

                    \begin{figure}[H]
                        \centering
                        \includegraphics[width=0.5\textwidth]{Problem/AHSME_I_1995_8_Problem_Figure_1.png}
                    \end{figure}
                    
$\mathrm{(A) \ 5 } \qquad \mathrm{(B) \ \frac {16}{3} } \qquad \mathrm{(C) \ \frac {20}{3} } \qquad \mathrm{(D) \ \frac {15}{2} } \qquad \mathrm{(E) \ 8 }$


 

	\item (\textit{AHSME 1995}) Consider the figure consisting of a square, its diagonals, and the segments joining the midpoints of opposite sides. The total number of triangles of any size in the figure is 

                    \begin{figure}[H]
                        \centering
                        \includegraphics[width=0.5\textwidth]{Problem/AHSME_I_1995_9_Problem_Figure_1.png}
                    \end{figure}
                    
$\mathrm{(A) \ 10 } \qquad \mathrm{(B) \ 12 } \qquad \mathrm{(C) \ 14 } \qquad \mathrm{(D) \ 16 } \qquad \mathrm{(E) \ 18 }$


 

	\item (\textit{AHSME 1995}) The area of the triangle bounded by the lines $y = x, y = - x$ and $y = 6$ is


 $\mathrm{(A) \ 12 } \qquad \mathrm{(B) \ 12\sqrt{2} } \qquad \mathrm{(C) \ 24 } \qquad \mathrm{(D) \ 24\sqrt{2} } \qquad \mathrm{(E) \ 36 }$


 

	\item (\textit{AHSME 1995}) How many base 10 four-digit numbers, $N = \underline{a} \underline{b} \underline{c} \underline{d}$, satisfy all three of the following conditions?


 $\text{(i)}$ $4,000 \leq N < 6,000;$


 $\text{(ii)}$ $N$ $\text{is a multiple of 5}$; 


 $\text{(iii)}$ $3 \leq b < c \leq 6$.


 

$\mathrm{(A) \ 10 } \qquad \mathrm{(B) \ 18 } \qquad \mathrm{(C) \ 24 } \qquad \mathrm{(D) \ 36 } \qquad \mathrm{(E) \ 48 }$


 

	\item (\textit{AHSME 1995}) Let $f$ be a linear function with the properties that $f(1) \leq f(2), f(3) \geq f(4),$ and $f(5) = 5$. Which of the following is true?


 $\mathrm{(A) \ f(0) < 0 } \qquad \mathrm{(B) \ f(0) = 0 } \qquad \mathrm{(C) \ f(1) < f(0) < f( - 1) } \qquad \mathrm{(D) \ f(0) = 5 } \qquad \mathrm{(E) \ f(0) > 5 }$


 

	\item (\textit{AHSME 1995}) The addition below is incorrect. The display can be made correct by changing one digit $d$, wherever it occurs, to another digit $e$. Find the sum of $d$ and $e$.


 $\begin{tabular}{ccccccc} & 7 & 4 & 2 & 5 & 8 & 6 \\ + & 8 & 2 & 9 & 4 & 3 & 0 \\ \hline 1 & 2 & 1 & 2 & 0 & 1 & 6 \end{tabular}$


 $\mathrm{(A) \ 4 } \qquad \mathrm{(B) \ 6 } \qquad \mathrm{(C) \ 8 } \qquad \mathrm{(D) \ 10 } \qquad \mathrm{(E) \ \text{more than 10} }$


 

	\item (\textit{AHSME 1995}) If $f(x) = ax^4 - bx^2 + x + 5$ and $f( - 3) = 2$, then $f(3) =$


 $\mathrm{(A) \ -5 } \qquad \mathrm{(B) \ -2 } \qquad \mathrm{(C) \ 1 } \qquad \mathrm{(D) \ 3 } \qquad \mathrm{(E) \ 8 }$


 

	\item (\textit{AHSME 1995}) Five points on a circle are numbered 1,2,3,4, and 5 in clockwise order. A bug jumps in a clockwise direction from one point to another around the circle; if it is on an odd-numbered point, it moves one point, and if it is on an even-numbered point, it moves two points. If the bug begins on point 5, after 1995 jumps it will be on point 

                    \begin{figure}[H]
                        \centering
                        \includegraphics[width=0.5\textwidth]{Problem/AHSME_I_1995_15_Problem_Figure_1.png}
                    \end{figure}
                    
$\mathrm{(A) \ 1 } \qquad \mathrm{(B) \ 2 } \qquad \mathrm{(C) \ 3 } \qquad \mathrm{(D) \ 4 } \qquad \mathrm{(E) \ 5 }$


 

	\item (\textit{AHSME 1995}) Anita attends a baseball game in Atlanta and estimates that there are 50,000 fans in attendance. Bob attends a baseball game in Boston and estimates that there are 60,000 fans in attendance. A league official who knows the actual numbers attending the two games notes that:


 i. The actual attendance in Atlanta is within $10 \%$ of Anita's estimate.
ii. Bob's estimate is within $10 \%$ of the actual attendance in Boston.


 To the nearest 1,000, the largest possible difference between the numbers attending the two games is


 $\mathrm{(A) \ 10000 } \qquad \mathrm{(B) \ 11000 } \qquad \mathrm{(C) \ 20000 } \qquad \mathrm{(D) \ 21000 } \qquad \mathrm{(E) \ 22000 }$


 

	\item (\textit{AHSME 1995}) Given regular pentagon $ABCDE$, a circle can be drawn that is tangent to $\overline{DC}$ at $D$ and to $\overline{AB}$ at $A$. The number of degrees in minor arc $AD$ is


 
                    \begin{figure}[H]
                        \centering
                        \includegraphics[width=0.5\textwidth]{Problem/AHSME_I_1995_17_Problem_Figure_1.png}
                    \end{figure}
                    


 $\mathrm{(A) \ 72 } \qquad \mathrm{(B) \ 108 } \qquad \mathrm{(C) \ 120 } \qquad \mathrm{(D) \ 135 } \qquad \mathrm{(E) \ 144 }$


 

	\item (\textit{AHSME 1995}) Two rays with common endpoint $O$ form a $30^\circ$ angle. Point $A$ lies on one ray, point $B$ on the other ray, and $AB = 1$. The maximum possible length of $OB$ is


 $\mathrm{(A) \ 1 } \qquad \mathrm{(B) \ \frac {1 + \sqrt {3}}{\sqrt 2} } \qquad \mathrm{(C) \ \sqrt{3} } \qquad \mathrm{(D) \ 2 } \qquad \mathrm{(E) \ \frac{4}{\sqrt{3}} }$


 

	\item (\textit{AHSME 1995}) Equilateral triangle $DEF$ is inscribed in equilateral triangle $ABC$ such that $\overline{DE} \perp \overline{BC}$. The ratio of the area of $\triangle DEF$ to the area of $\triangle ABC$ is

                    \begin{figure}[H]
                        \centering
                        \includegraphics[width=0.5\textwidth]{Problem/AHSME_I_1995_19_Problem_Figure_1.png}
                    \end{figure}
                    
$\mathrm{(A) \ \frac {1}{6} } \qquad \mathrm{(B) \ \frac {1}{4} } \qquad \mathrm{(C) \ \frac {1}{3} } \qquad \mathrm{(D) \ \frac {2}{5} } \qquad \mathrm{(E) \ \frac {1}{2} }$


 

	\item (\textit{AHSME 1995}) If $a,b$ and $c$ are three (not necessarily different) numbers chosen randomly and with replacement from the set $\{1,2,3,4,5 \}$, the probability that $ab + c$ is even is


 $\mathrm{(A) \ \frac {2}{5} } \qquad \mathrm{(B) \ \frac {59}{125} } \qquad \mathrm{(C) \ \frac {1}{2} } \qquad \mathrm{(D) \ \frac {64}{125} } \qquad \mathrm{(E) \ \frac {3}{5} }$


 

	\item (\textit{AHSME 1995}) Two nonadjacent vertices of a rectangle are $(4,3)$ and $(-4,-3)$, and the coordinates of the other two vertices are integers. The number of such rectangles is


 $\mathrm{(A) \ 1 } \qquad \mathrm{(B) \ 2 } \qquad \mathrm{(C) \ 3 } \qquad \mathrm{(D) \ 4 } \qquad \mathrm{(E) \ 5 }$


 

	\item (\textit{AHSME 1995}) A pentagon is formed by cutting a triangular corner from a rectangular piece of paper. The five sides of the pentagon have lengths 13, 19, 20, 25 and 31, although this is not necessarily their order around the pentagon. The area of the pentagon is


 $\mathrm{(A) \ 459 } \qquad \mathrm{(B) \ 600 } \qquad \mathrm{(C) \ 680 } \qquad \mathrm{(D) \ 720 } \qquad \mathrm{(E) \ 745 }$


 

	\item (\textit{AHSME 1995}) The sides of a triangle have lengths 11, 15, and $k$, where $k$ is an integer. For how many values of $k$ is the triangle obtuse?


 $\mathrm{(A) \ 5 } \qquad \mathrm{(B) \ 7 } \qquad \mathrm{(C) \ 12 } \qquad \mathrm{(D) \ 13 } \qquad \mathrm{(E) \ 14 }$


 

	\item (\textit{AHSME 1995}) There exist positive integers $A,B$ and $C$, with no common factor greater than 1, such that


 
 \[A \log_{200} 5 + B \log_{200} 2 = C\]


 What is $A + B + C$?


 $\mathrm{(A) \ 6 } \qquad \mathrm{(B) \ 7 } \qquad \mathrm{(C) \ 8 } \qquad \mathrm{(D) \ 9 } \qquad \mathrm{(E) \ 10 }$


 

	\item (\textit{AHSME 1995}) A list of five positive integers has mean 12 and range 18. The mode and median are both 8. How many different values are possible for the second- largest element of the list? 


 $\mathrm{(A) \ 4 } \qquad \mathrm{(B) \ 6 } \qquad \mathrm{(C) \ 8 } \qquad \mathrm{(D) \ 10 } \qquad \mathrm{(E) \ 12 }$


 

	\item (\textit{AHSME 1995}) In the figure, $\overline{AB}$ and $\overline{CD}$ are diameters of the circle with center $O$, $\overline{AB} \perp \overline{CD}$, and chord $\overline{DF}$ intersects $\overline{AB}$ at $E$. If $DE = 6$ and $EF = 2$, then the area of the circle is 


 
                    \begin{figure}[H]
                        \centering
                        \includegraphics[width=0.5\textwidth]{Problem/AHSME_I_1995_26_Problem_Figure_1.png}
                    \end{figure}
                    


 $\mathrm{(A) \ 23 \pi } \qquad \mathrm{(B) \ \frac {47}{2} \pi } \qquad \mathrm{(C) \ 24 \pi } \qquad \mathrm{(D) \ \frac {49}{2} \pi } \qquad \mathrm{(E) \ 25 \pi }$


 

	\item (\textit{AHSME 1995}) Consider the triangular array of numbers with 0,1,2,3,... along the sides and interior numbers obtained by adding the two adjacent numbers in the previous row. Rows 1 through 6 are shown.


 
 \[\begin{tabular}{ccccccccccc} & & & & & 0 & & & & & \\ & & & & 1 & & 1 & & & & \\ & & & 2 & & 2 & & 2 & & & \\ & & 3 & & 4 & & 4 & & 3 & & \\ & 4 & & 7 & & 8 & & 7 & & 4 & \\ 5 & & 11 & & 15 & & 15 & & 11 & & 5  \end{tabular}\]


 Let $f(n)$ denote the sum of the numbers in row $n$. What is the remainder when $f(100)$ is divided by 100?


 $\mathrm{(A) \ 12 } \qquad \mathrm{(B) \ 30 } \qquad \mathrm{(C) \ 50 } \qquad \mathrm{(D) \ 62 } \qquad \mathrm{(E) \ 74 }$


 

	\item (\textit{AHSME 1995}) Two parallel chords in a circle have lengths 10 and 14, and the distance between them is 6. The chord parallel to these chords and midway between them is of length $\sqrt {a}$ where $a$ is

                    \begin{figure}[H]
                        \centering
                        \includegraphics[width=0.5\textwidth]{Problem/AHSME_I_1995_28_Problem_Figure_1.png}
                    \end{figure}
                    
$\mathrm{(A) \ 144 } \qquad \mathrm{(B) \ 156 } \qquad \mathrm{(C) \ 168 } \qquad \mathrm{(D) \ 176 } \qquad \mathrm{(E) \ 184 }$


 

	\item (\textit{AHSME 1995}) For how many three-element sets of distinct positive integers $\{a,b,c\}$ is it true that $a \times b \times c = 2310$?


 $\mathrm{(A) \ 32 } \qquad \mathrm{(B) \ 36 } \qquad \mathrm{(C) \ 40 } \qquad \mathrm{(D) \ 43 } \qquad \mathrm{(E) \ 45 }$


 

	\item (\textit{AHSME 1995}) A large cube is formed by stacking 27 unit cubes. A plane is perpendicular to one of the internal diagonals of the large cube and bisects that diagonal. The number of unit cubes that the plane intersects is

                    \begin{figure}[H]
                        \centering
                        \includegraphics[width=0.5\textwidth]{Problem/AHSME_I_1995_30_Problem_Figure_1.png}
                    \end{figure}
                    
$\mathrm{(A) \ 16 } \qquad \mathrm{(B) \ 17 } \qquad \mathrm{(C) \ 18 } \qquad \mathrm{(D) \ 19 } \qquad \mathrm{(E) \ 20 }$


 

\end{enumerate}